\section{Bewertung und priorisierte Maßnahmen}
\label{sec:massnahmen}

\subsection{Ergebnisübersicht}

\begin{table}[H]
  \centering
  \caption{Bewertung zum Berichtsstand nach XSS-Rollout}
  \small
  \begin{tabularx}{\textwidth}{l l X}
    \toprule
    \textbf{Bereich} & \textbf{Status} & \textbf{Begründung}\\
    \midrule
    VueCal Stored DOM XSS & behoben & Eigene Event-Slots und
    Vue-Interpolation sind auf \texttt{capstone-main} vollständig
    ausgerollt; erneute Prüfung ohne aktuellen Befund.\\
    Eigene Vue-HTML-Senken & kein Befund & Keine produktive direkte
    Verwendung von \texttt{v-html}, \texttt{innerHTML} oder
    vergleichbaren Senken in \texttt{web/src}.\\
    Spring API-CORS & kein Allow-All-Befund & Ohne Demo-Konfiguration
    keine explizite globale CORS-Freigabe; normale Entwicklung nutzt
    den Vite-Proxy.\\
    \texttt{/files/**} & hoch & Voll-Exporte und Token-Erzeugung
    liegen unter einer globalen \texttt{permitAll}-Regel.\\
    Keycloak Web-Client & hoch & Wildcard-Web-Origin und nicht
    benötigte beziehungsweise veraltete OAuth-Flows.\\
    Security-Regression & ausbaufähig & Spezifische automatisierte
    XSS-, CORS- und Autorisierungstests fehlen im regulären Stand.\\
    \bottomrule
  \end{tabularx}
\end{table}

\subsection{Priorität 1: Öffentliche Dateioberfläche begrenzen}

Die pauschale Freigabe von \texttt{/files/**} sollte durch getrennte
Regeln ersetzt werden. Nur ausdrücklich öffentliche, ausreichend
zufällige und möglichst zeitlich begrenzte Token-Downloads dürfen
ohne Login erreichbar sein. Folgende Endpunkte müssen mindestens
authentifiziert und fachlich autorisiert werden:

\begin{itemize}
  \item \texttt{GET /files/planningperiod/\{id\}/download},
  \item sämtliche \texttt{POST /files/create/...}-Endpunkte,
  \item \texttt{POST /files/get-or-create-ics-token}.
\end{itemize}

\subsection{Priorität 2: XSS-Fix durch Tests absichern}

Der ausgerollte Fix sollte mit Komponenten- oder End-to-End-Tests
gegen Regression geschützt werden. Für jede \texttt{VueCal}-Ansicht
wird ein Event mit HTML- und Event-Handler-Zeichenfolge gerendert.
Der Test muss bestätigen, dass der String als Text sichtbar bleibt
und kein DOM-Element aus dem Payload entsteht. Ergänzend sollte CI
nach neuen Verwendungen von \texttt{v-html}, \texttt{innerHTML} und
ähnlichen Senken suchen.

\subsection{Priorität 3: Keycloak und OAuth härten}

In \path{scheduler.json} sind konkrete Frontend-Origins statt
\texttt{*} zu hinterlegen. Implicit Flow und Direct Access Grants
sollten für den SPA-Web-Client deaktiviert, der Authorization Code
Flow mit PKCE verwendet und der Brute-Force-Schutz aktiviert werden.
Außerhalb lokaler Entwicklung muss TLS erzwungen werden.

\subsection{Priorität 4: CORS nur bei fachlichem Bedarf aktivieren}

Solange Frontend und API über denselben Origin beziehungsweise
Reverse Proxy erreichbar sind, ist keine globale API-CORS-Freigabe
nötig. Falls Cross-Origin-Zugriff eingeführt wird, sollte
\texttt{SecurityConfiguration.java} eine properties-basierte
Allowlist mit blockierendem Default erhalten. Die Registrierung
sollte auf tatsächlich benötigte Routen begrenzt werden. Demo-Origins
wie \texttt{localhost:5173} dürfen ausschließlich in einem lokalen
Entwicklungsprofil vorkommen.

\subsection{Priorität 5: Zusätzliche Schutzschichten}

\begin{itemize}
  \item \texttt{credentials: 'include'} in
    \texttt{web/src/shared/http-client.js} auf \texttt{same-origin}
    oder \texttt{omit} reduzieren.
  \item CSP im produktiven Reverse Proxy oder Backend einführen und
    zunächst im Report-Only-Modus testen.
  \item Freitextfelder serverseitig auf fachlich sinnvolle Längen und
    Werte validieren.
  \item \texttt{window.\$keycloak} entfernen oder den globalen
    Zugriff mindestens reduzieren.
  \item Dependency- und Security-Checks für npm und Gradle in CI integrieren.
  \item Die produktive Nginx-/Ingress-Konfiguration dokumentieren;
    die eingecheckte \texttt{nginx.conf} enthält derzeit keine
    \texttt{/api}-Proxy-Regel.
\end{itemize}
